% GENERATED FILE. Edit canonical lessons, then run npm run generate:books.
% canonical-source-hash: fnv1a64:5b8e50c8029b0eec
% canonical-lessons: AR-C496-mutahammis, AR-C496-tahammasa, AR-C496-irtiyah, AR-C496-murtah, AR-C496-irtaha

\chapter{Excited and relieved}
\label{ch:ar-excited-and-relieved}
\hlchaptermodality{496}

\begin{chapteropening}
\textbf{By the end of this chapter:} \emph{I can say that someone is excited or at ease, talk about relief, and say that someone got excited or felt relieved.}

Complete the last lesson of chapter 496: I can say that someone is excited or at ease, talk about relief, and say that someone got excited or felt relieved.
\end{chapteropening}

\section[\texorpdfstring{\ar{متحمس} (mutaḥammis)}{متحمس (mutaḥammis)}]{\ar{متحمس} (mutaḥammis) — excited, keen}

\label{lesson:AR-C496-mutahammis}



\begin{quote}
Before the new one: say the Arabic for to put up with, then the Arabic for to be fed up.
\end{quote}

\subsection*{You'll want to know: \ar{متحمس}}
\textbf{\ar{متحمس}} — \emph{mutaḥammis} — \textquotedblleft{}excited, keen\textquotedblright{}.

\begin{quote}
\ar{ابني} \ar{متحمس} \ar{للرحلة} \ar{إلى} \ar{البحر}.
\end{quote}

\begin{quote}
\emph{ibnī mutaḥammis li-r-riḥla ilā al-baḥr.}
\end{quote}

\begin{quote}
My son is excited about the trip to the sea.
\end{quote}

It takes \ar{لـ} for what you look forward to.

\begin{cousinweb}
\textbf{\ar{متحمس}} (\emph{mutaḥammis}) is built on the native Arabic root \emph{ḥ-m-s}, \textquotedblleft{}to be hard, firm, vehement\textquotedblright{} (Lane). Ḥamās, enthusiasm, which you know, is from the same root.
\end{cousinweb}

\begin{grammarlens}[title={What you've built}]
One more word for describing what you lived and felt.
\end{grammarlens}

\subsection*{Your turn}
\begin{itemize}
\raggedright
  \item \emph{Say it:} \emph{mutaḥammis}
  \item \emph{Say it:} \emph{mutaḥammis}, once more
  \item \emph{Say it:} \emph{ibnī mutaḥammis li-r-riḥla ilā al-baḥr.} — My son is excited about the trip to the sea.
  \item \emph{Recall:} say the Arabic for to put up with, then the Arabic for to be fed up, then say \emph{mutaḥammis} again
\end{itemize}

\begin{tcolorbox}[breakable,colback=teal!4,colframe=teal!35!black,title={Before you move on}]
What does \ar{متحمس} mean? (\textquotedblleft{}Excited, keen\textquotedblright{}.) Say it once more.
\end{tcolorbox}
\section[\texorpdfstring{\ar{تحمس} (taḥammasa)}{تحمس (taḥammasa)}]{\ar{تحمس} (taḥammasa) — to get excited}

\label{lesson:AR-C496-tahammasa}



\begin{quote}
Before the new one: say the Arabic for to be fed up, then the Arabic for excited.
\end{quote}

\subsection*{You'll want to know: \ar{تحمس}}
\textbf{\ar{تحمس}} — \emph{taḥammasa} — \textquotedblleft{}to get excited\textquotedblright{}.

\begin{quote}
\ar{تحمس} \ar{الأولاد} \ar{عندما} \ar{بدأت} \ar{المباراة}.
\end{quote}

\begin{quote}
\emph{taḥammasa al-awlād ʿindamā badaʾat al-mubārā.}
\end{quote}

\begin{quote}
The boys got excited when the match began.
\end{quote}

Present: \ar{يتحمس}, \emph{yataḥammas}.

\begin{cousinweb}
\textbf{\ar{تحمس}} (\emph{taḥammasa}) is built on the native Arabic root \emph{ḥ-m-s}, \textquotedblleft{}to be hard, firm, vehement\textquotedblright{} (Lane). Form V taḥammasa is to grow keen.
\end{cousinweb}

\begin{grammarlens}[title={What you've built}]
One more word for describing what you lived and felt.
\end{grammarlens}

\subsection*{Your turn}
\begin{itemize}
\raggedright
  \item \emph{Say it:} \emph{taḥammasa}
  \item \emph{Say it:} \emph{taḥammasa}, once more
  \item \emph{Say it:} \emph{taḥammasa al-awlād ʿindamā badaʾat al-mubārā.} — The boys got excited when the match began.
  \item \emph{Recall:} say the Arabic for to be fed up, then the Arabic for excited, then say \emph{taḥammasa} again
\end{itemize}

\begin{tcolorbox}[breakable,colback=teal!4,colframe=teal!35!black,title={Before you move on}]
What does \ar{تحمس} mean? (\textquotedblleft{}To get excited\textquotedblright{}.) Say it once more.
\end{tcolorbox}
\section[\texorpdfstring{\ar{ارتياح} (irtiyāḥ)}{ارتياح (irtiyāḥ)}]{\ar{ارتياح} (irtiyāḥ) — relief}

\label{lesson:AR-C496-irtiyah}



\begin{quote}
Before the new one: say the Arabic for excited, then the Arabic for to get excited.
\end{quote}

\subsection*{You'll want to know: \ar{ارتياح}}
\textbf{\ar{ارتياح}} — \emph{irtiyāḥ} — \textquotedblleft{}relief\textquotedblright{}.

\begin{quote}
\ar{كان} \ar{الارتياح} \ar{واضحا} \ar{على} \ar{وجهها}.
\end{quote}

\begin{quote}
\emph{kāna al-irtiyāḥ wāḍiḥan ʿalā wajhihā.}
\end{quote}

\begin{quote}
Relief was clear on her face.
\end{quote}

\begin{cousinweb}
\textbf{\ar{ارتياح}} (\emph{irtiyāḥ}) is built on the native Arabic root \emph{r-w-ḥ}, \textquotedblleft{}wind, breath, rest\textquotedblright{} (Lane). Irtiyāḥ is the verbal noun of Form VIII: breathing easy again.
\end{cousinweb}

\begin{grammarlens}[title={What you've built}]
One more word for describing what you lived and felt.
\end{grammarlens}

\subsection*{Your turn}
\begin{itemize}
\raggedright
  \item \emph{Say it:} \emph{irtiyāḥ}
  \item \emph{Say it:} \emph{irtiyāḥ}, once more
  \item \emph{Say it:} \emph{kāna al-irtiyāḥ wāḍiḥan ʿalā wajhihā.} — Relief was clear on her face.
  \item \emph{Recall:} say the Arabic for excited, then the Arabic for to get excited, then say \emph{irtiyāḥ} again
\end{itemize}

\begin{tcolorbox}[breakable,colback=teal!4,colframe=teal!35!black,title={Before you move on}]
What does \ar{ارتياح} mean? (\textquotedblleft{}Relief\textquotedblright{}.) Say it once more.
\end{tcolorbox}
\section[\texorpdfstring{\ar{مرتاح} (murtāḥ)}{مرتاح (murtāḥ)}]{\ar{مرتاح} (murtāḥ) — at ease, relaxed}

\label{lesson:AR-C496-murtah}



\begin{quote}
Before the new one: say the Arabic for to get excited, then the Arabic for relief.
\end{quote}

\subsection*{You'll want to know: \ar{مرتاح}}
\textbf{\ar{مرتاح}} — \emph{murtāḥ} — \textquotedblleft{}at ease, relaxed\textquotedblright{}.

\begin{quote}
\ar{أنا} \ar{مرتاح} \ar{في} \ar{عملي} \ar{الجديد}.
\end{quote}

\begin{quote}
\emph{anā murtāḥ fī ʿamalī al-jadīd.}
\end{quote}

\begin{quote}
I am at ease in my new job.
\end{quote}

\begin{cousinweb}
\textbf{\ar{مرتاح}} (\emph{murtāḥ}) is built on the native Arabic root \emph{r-w-ḥ}, \textquotedblleft{}wind, breath, rest\textquotedblright{} (Lane). Murtāḥ is the participle of Form VIII.
\end{cousinweb}

\begin{grammarlens}[title={What you've built}]
One more word for describing what you lived and felt.
\end{grammarlens}

\subsection*{Your turn}
\begin{itemize}
\raggedright
  \item \emph{Say it:} \emph{murtāḥ}
  \item \emph{Say it:} \emph{murtāḥ}, once more
  \item \emph{Say it:} \emph{anā murtāḥ fī ʿamalī al-jadīd.} — I am at ease in my new job.
  \item \emph{Recall:} say the Arabic for to get excited, then the Arabic for relief, then say \emph{murtāḥ} again
\end{itemize}

\begin{tcolorbox}[breakable,colback=teal!4,colframe=teal!35!black,title={Before you move on}]
What does \ar{مرتاح} mean? (\textquotedblleft{}At ease, relaxed\textquotedblright{}.) Say it once more.
\end{tcolorbox}
\section[\texorpdfstring{\ar{ارتاح} (irtāḥa)}{ارتاح (irtāḥa)}]{\ar{ارتاح} (irtāḥa) — to feel relieved}

\label{lesson:AR-C496-irtaha}



\begin{quote}
Before the new one: say the Arabic for relief, then the Arabic for at ease.
\end{quote}

\subsection*{You'll want to know: \ar{ارتاح}}
\textbf{\ar{ارتاح}} — \emph{irtāḥa} — \textquotedblleft{}to feel relieved\textquotedblright{}.

\begin{quote}
\ar{ارتاح} \ar{أبي} \ar{عندما} \ar{وصلنا} \ar{إلى} \ar{البيت}.
\end{quote}

\begin{quote}
\emph{irtāḥa abī ʿindamā waṣalnā ilā al-bayt.}
\end{quote}

\begin{quote}
My father felt relieved when we reached home.
\end{quote}

Present: \ar{يرتاح}, \emph{yartāḥ}.

The same root gives \ar{استراح}, to rest, which you know.

\begin{cousinweb}
\textbf{\ar{ارتاح}} (\emph{irtāḥa}) is built on the native Arabic root \emph{r-w-ḥ}, \textquotedblleft{}wind, breath, rest\textquotedblright{} (Lane). It is Form VIII of the root.
\end{cousinweb}

\begin{grammarlens}[title={What you've built}]
One more word for describing what you lived and felt.
\end{grammarlens}

\subsection*{Your turn}
\begin{itemize}
\raggedright
  \item \emph{Say it:} \emph{irtāḥa}
  \item \emph{Say it:} \emph{irtāḥa}, once more
  \item \emph{Say it:} \emph{irtāḥa abī ʿindamā waṣalnā ilā al-bayt.} — My father felt relieved when we reached home.
  \item \emph{Recall:} say the Arabic for relief, then the Arabic for at ease, then say \emph{irtāḥa} again
\end{itemize}

\begin{tcolorbox}[breakable,colback=teal!4,colframe=teal!35!black,title={Before you move on}]
What does \ar{ارتاح} mean? (\textquotedblleft{}To feel relieved\textquotedblright{}.) Say it once more.
\end{tcolorbox}
